We now provide adaptive concentration bounds for the random cumulative incidence forests. Note that the following shows consistency over time $t$ only, not over $\histlower \in \history$.
Let $\scriptNode = \{ \Nodeu \}_{u \in U}$ be the set of terminal nodes obtained by fitting a single tree, where $U$ is a set of indices, $Z = \bigcup_{u \in U} \Nodeu$, and $\Nodeu \cap \Node_l = \emptyset$ for any $u \neq l$. Recursive partitioning is a critical component of tree-based survival models, where a $d$-dimensional feature space $Z$ is partitioned into terminal nodes. In other words, $\scriptNode$ is a partition of $Z$. Following the traditional tree-building process in \cite{breiman1984classification}, binary splitting rules are used, and all terminal nodes $\Nodeu$ are non-overlapping hyper-rectangles \citep{wager_adaptive_2016, cui_consistency_2022, cho_multi-stage_2023}. A set of $B$ trees are fitted in a traditional survival forest model \citep{ishwaran_random_2008}. 
Hence, our tree model yields $\{ \Nodeu, \hat{\cuminc}_{\Nodeu,n}\}_{u \in U}$, and our forest constructs a collection of partitions $\{\{ \Nodeu^b, \hat{\cuminc}_{\Nodeu,n} \}_{u \in U_b}\}^B_{b=1}$.

The AJ estimator for any node $\Nodeu$ is
\begin{equation}
\label{eqn:AJhat_node}
\hat{\cuminc}_{\Nodeu,n}(t) = 
\sum_{s \leq t}
\frac{
\survfunchatNodeu(s-)
    \sum_{i=1}^n I{\{ \delta_{1i} = 1, X_i = s  \}}
    I{\{ Z_i \in \Nodeu \}}}{\sum_{i=1}^n I{\{ X_i \geq s \}} I{\{ Z_i \in \Nodeu \}}},
\end{equation} where the Kaplan-Meier estimator for overall survival in node $\Nodeu$ at time $t$ is given by 
$\survfunchatNodeu(t)$.
We know that the Kaplan-Meier estimator in the survival forest is consistent. To quantify within-node estimation performance, let $\cuminc_\Node(t) = \frac{1}{\mu(\Node)} \int_{z\in \Node} \cuminc(t | z) dP(z)$ be the averaged CIF within any node $\Node$, where $\cuminc(t|z)$ is defined in Equation \cref{eqn:CIF_def}, $P$ is the distribution of $Z$, and $\mu(\Node) = \int_{z\in \Node} dP(z)$ is the measure of node $\mathcal{Z}$.

Equation \cref{eqn:AJhat_node} can be re-written in terms of distributions and hazard functions. We already know that $X_i = T_i \wedge C_i$ and $T_i = T_{1i} \wedge T_{2i}$. Let $H_i(t)$ be the CDF of $T_i$ and $G_i(t)$ be the CDF of $C_i$. 
%We also defined $N_i(s) = I(T_i \leq s, \epsilon_i = 1)$ where we assumed cause 1 is the priority cause, and $Y_i(s) = 1-N_i(s-)$ \citep{fine_proportional_1999}. \\ Then, $dN_i(s) = I(T_i = s, \epsilon_i = 1)$.
The counting process for cause-$j$ is given by $N_j(s) = \sum_{i=1}^{n} N_{ji}(s) = \sum_{i=1}^{n} 1(X_i \leq s, \, \delta_{ji} = 1)$, and the at-risk process is given by
$ K(s) = \sum_{i=1}^{n} K_i(s) = \sum_{i=1}^{n} 1(X_i \geq s)$. Then, taking the expectation of the risk set, we get the denominator to be
\begin{equation}
\begin{split}
\mathbb{E}(K_i(s)) &= \mathbb{E}(I(X_i \geq s)) = \Pr(X_i \geq s) = \Pr(T_i \wedge C_i \geq s) \\
&= \Pr(T_i \geq s, C_i \geq s) = \Pr(T_i \geq s) \Pr(C_i \geq s) \\
&=[1-\Pr(T_i < s)][1 - \Pr(C_i < s)]\\
&= [1-H_i(s)][1-G_i(s)],
\end{split}
\end{equation}
where we assume independence between overall failure and censoring times.
Next, we take the expectation of the counts of failures from cause $j$ at time $s$ to get the numerator:
\begin{equation}
    \begin{split}
\mathbb{E}(dN_{ji}(s)) &= \mathbb{E}(I(\delta_{ji} = 1, X_i = s)) \\ 
&= \mathbb{E}(I(\delta_{ji} = 1, T_i \wedge C_i = s, \epsilon_i = j)) \\ 
&= \Pr(\delta_{ji} = 1, T_i \wedge C_i = s, \epsilon_i = j, X_i \geq s) \\
&= \Pr(\delta_{ji} = 1, T_i \wedge C_i = s, \epsilon_i = j | X_i \geq s) \Pr(X_i > s)\\
&= \Pr(\delta_{ji} = 1, T_i \wedge C_i = s, \epsilon_i = j | T_i \geq s, C_i \geq s) \Pr(T_i > s, C_i > s)\\
&= \Pr(\delta_{ji} = 1, T_i \wedge C_i = s, \epsilon_i = j | T_i \geq s) [1-\Pr(T_i \leq s)][1 - \Pr(C_i \leq s)]\\
&= d\Lambda_{ji}(s)[1-H_i(s)][1-G_i(s)],
\end{split}
\end{equation}
where $d\Lambda_j(s)$ is the cause-$j$ specific hazard function defined in Equation \cref{eqn:dLambdaj}. The third equality holds because $X_i \geq s$ is a larger set than $X_i = s$, and the sixth equality holds because we assume independent censoring (i.e., being conditional on $T_i,C_i$ is the same as being conditional on just $T_i$ because $T_i \perp C_i$).


